\documentclass[10pt,a4paper]{amsart}
\usepackage{amsmath,amssymb,amsthm,graphicx}
\usepackage{url}
\usepackage{hyperref}
\usepackage{enumerate}

\DeclareMathOperator{\Cliff}{Clif\/f}
\DeclareMathOperator{\Gal}{Gal}
\DeclareMathOperator{\GL}{GL}
\DeclareMathOperator{\Inn}{Inn}
\DeclareMathOperator{\Mat}{Mat}
\DeclareMathOperator{\Out}{Out}
\DeclareMathOperator{\SL}{SL}
\DeclareMathOperator{\SO}{SO}
\DeclareMathOperator{\SU}{\mathbf{SU}}
\numberwithin{equation}{section}

\newtheorem{prop}{Proposition}[section]
\newtheorem{cor}[prop]{Corollary}
\newtheorem{lem}[prop]{Lemma}
{\theoremstyle{definition}
\newtheorem{defn}[prop]{Definition}
\newtheorem{notation}[prop]{Notation}
\newtheorem{rem}[prop]{Remark}
\newtheorem{ques}[prop]{Question}
\newtheorem{eg}[prop]{Example}
\newtheorem{erratum}[prop]{Erratum}}
\newcommand{\Res}[1]{\mathop{\mathbf{R}_{#1}}}

\newcommand{\ResOne}[1]{\mathop{\mathbf{R}^{(1)}_{#1}}}

\newcommand{\Lie}[1]{\mathfrak{\lowercase{#1}}}

\newcommand{\arXivNumber}{1410.2339}

\allowdisplaybreaks
\title{Surjectivity of the restriction map in Galois cohomology from an imaginary quadratic extension to the real place for a connected semisimple algebraic group over the rationals}
\author{Dave Witte Morris}
\date{doi:10.3842/sigma.2015.034; CC BY 4.0}
\begin{document}
\maketitle
\setcounter{section}{2}
\setcounter{prop}{4}
\noindent\textit{Editorial note.} The counted unit is the author's labelled lemma KneserHasse at source lines 245-249 with the proof environment that begins two lines later at source line 251; the pairing is the ordinary adjacent one. The optional name of the lemma carries the citation to Kneser's lectures from which the statement is taken, whereas the proof printed here is the author's own argument, credited in the proof itself to A. Rapinchuk.\par\medskip
\noindent\textit{Editorial note.} The retained counter values make the counted lemma print as 2.5, its number in the pinned source: the statement is the first numbered item of the excerpt and four proposition-class items precede it in source section 2.\par\medskip
\noindent\textit{Editorial note.} The retained preamble is the author's own package, operator and theorem block followed by the author's own macros, taken from the pinned native source; the publisher sigma document class is not part of the retained material, and no mathematical macro, symbol, statement or proof step was added, changed or removed.\par\medskip
\noindent\textit{Editorial note.} The retained bibliography is the article's own reference list with the publisher's trailing page-footer command omitted; the copy kept with this task exists only so that the material is self-contained, and no entry was rewritten.\par\medskip
\noindent\textit{Editorial note.} No retained line refers to a label outside the retained spans, so the package declares no external reference.\par\medskip
\noindent\textit{Allowed standard prerequisites (editorial).}\begin{itemize}
\item Galois cohomology of linear algebraic groups over a field of characteristic zero. Only the standard formalism of first Galois cohomology sets, restriction of scalars and the identification of the relevant classes of a torus with the elements satisfying t t\textasciicircum{}sigma = 1 is assumed. The construction that produces a representative defined over an imaginary quadratic field is the retained obligation.
\item Weak approximation for tori at the infinite place. The density of the rational points of a rational torus in its real points is quoted as a standard arithmetic fact; every step that transports the cohomology class along that density is retained.
\item Restriction of scalars for tori and the norm-one subtorus. The torus obtained from T by restriction of scalars from L to Q, the Q-automorphism tau and the subtorus of elements satisfying r r\textasciicircum{}tau = 1 are constructed inside the retained proof; nothing about them is assumed beyond the standard properties of restriction of scalars.
\end{itemize}\medskip
We begin by recalling a~result of J.~Tits that uses Galois cohomology to characterize the irreducible representations of
semisimple algebraic groups over f\/ields that are not algebraically closed.
However, we will state this work in the setting of Lie algebras, rather than algebraic groups.
(We deal only with semisimple Lie algebras and semisimple groups, and the f\/ields under consideration in this paper are
all of characteristic zero, so no dif\/f\/iculties arise in making this translation~\cite[Proposition~7.3.1(iii), p.~393]{DemazureGrothendieck-SGA3}.)

\begin{lem}[\protect{\cite[Theorem~5.1b, p.~77]{Kneser-LectGalCohoClassical}}]\label{KneserHasse}
If $\mathbf{G}$ is a~connected, semisimple algebraic group over~$\mathbb{Q}$, and~$L$ is any imaginary quadratic
extension of~$\mathbb{Q}$, then the natural restriction map $H^1 \bigl(L/F; \mathbf{G}(L) \bigr) \to H^1 (\mathbb{R};
\mathbf{G})$ is surjective.
\end{lem}

\begin{proof}(The author thanks A.~Rapinchuk for suggesting this argument.) Let~$\sigma$ be the nontrivial element of~$\Gal(\mathbb{C}/\mathbb{R})$.
It is well known that any cohomology class in $H^1(\mathbb{R}; \mathbf{G})$ is represented by an element~$t$ of
a~maximal torus $\mathbf{T}(\mathbb{C})$ of $\mathbf{G}(\mathbb{C})$, such that $t t^\sigma = 1$, and, since the variety
of maximal tori has weak approximation, that $\mathbf{T}$ may be chosen to be def\/ined over~$\mathbb{Q}$ (cf.\ proof
of~\cite[Proposition~6.17, p.~337]{PlatonovRapinchukBook}).

Let $\mathbf{R} = \mathrm{Res}_{L/\mathbb{Q}} \mathbf{T}$ be the torus obtained from~$\mathbf{T}$ by restriction of
scalars from~$L$ to~$\mathbb{Q}$.
Then there is a~natural isomorphism $\varphi \colon \mathbf{T}(\mathbb{C}) \to \mathbf{R}(\mathbb{R})$, such that
$\varphi \bigl(\mathbf{T}(L) \bigr) = \mathbf{R}(\mathbb{Q})$, and there is a~$\mathbb{Q}$-automorphism~$\tau$
of~$\mathbf{R}$, such that $\varphi(x^\sigma) = \varphi(x)^\tau$, for all $x \in \mathbf{T}(\mathbb{C})$.

Let $\mathbf{R}^{(1)} = \{ r \in \mathbf{R} \mid r r^\tau = 1 \}$.
We claim that this is a~subtorus of~$\mathbf{R}$.
Since $\mathbf{T}$ is def\/ined over~$\mathbb{Q}$, there is an~$L$-isomorphism $\psi \colon \mathbf{R} \to \mathbf{T}
\times \mathbf{T}$, such that $\psi \bigl(\mathbf{R}(\mathbb{R}) \bigr) = \{ (x, x^\sigma) \mid x \in
\mathbf{T}(\mathbb{C}) \}$.
Therefore, if we def\/ine $(x,y)^\pi = (y,x)$, then $\psi(r^\tau) = \psi(r)^\pi$ for all $r \in \mathbf{R}$.
So $\psi(\mathbf{R}^{(1)}) = \{ (x, x^{-1}) \mid x \in \mathbf{T} \}$ is a~subtorus of $\mathbf{T} \times \mathbf{T}$,
as claimed.

Now, $\mathbf{R}^{(1)}$ is a~torus that is def\/ined over~$\mathbb{Q}$ (since the automorphism~$\tau$ is def\/ined
over~$\mathbb{Q}$), and all $\mathbb{Q}$-tori have weak approximation at the inf\/inite place~\cite[Corollary~1 of Proposition~7.8,
p.~418]{PlatonovRapinchukBook}, so $\mathbf{R}^{(1)}(\mathbb{Q})$ is dense in $\mathbf{R}^{(1)}(\mathbb{R})$.
Hence, some $q \in \mathbf{R}^{(1)}(\mathbb{Q})$ is in the same connected component of $\mathbf{R}^{(1)}(\mathbb{R})$ as
$\varphi(t)$.
Then, letting $q' = \varphi^{-1}(q)$, we see that $q'$ is in the same connected component of $\{ w \in
\mathbf{T}(\mathbb{R}) \mid w w^\sigma = 1 \}$ as~$t$, so $q'$ represents the same cohomology class as~$t$ in
$H^1(\mathbb{R}; \mathbf{G})$.
However, since $q' \in \mathbf{T}(L)$, we see that the cohomology class of~$q'$ is in the image of $H^1 \bigl(L/F;
\mathbf{G}(L) \bigr)$, as desired.
\end{proof}

\begin{thebibliography}{99}
\footnotesize\itemsep=0pt

\bibitem{DemazureGrothendieck-SGA3}
Demazure M., Grothendieck A., Sch\'emas en groups.~III. Structure des sch\'emas
  en groupes reductifs, S\'eminaire de g\'eom\'etrie alg\'ebrique du {B}ois
  {M}arie, 1962/1964, available at
  \url{http://library.msri.org/books/sga/sga/pdf/sga3-3.pdf}.

\bibitem{Harder-BerichtNeuere}
Harder G., Bericht \"uber neuere {R}esultate der {G}aloiskohomologie
  halbeinfacher {G}ruppen, \textit{Jber. Deutsch. Math.-Verein.} \textbf{70}
  (1968), 182--216, available at
  \url{http://www.digizeitschriften.de/en/dms/img/?PPN=PPN37721857X_0070&DMDID=dmdlog17}.

\bibitem{Kneser-LectGalCohoClassical}
Kneser M., Lectures on {G}alois cohomology of classical groups, \textit{Tata
  Institute of Fundamental Research Lectures on Mathematics}, Vol.~47, Tata
  Institute of Fundamental Research, Bombay, 1969, available at
  \url{http://www.math.tifr.res.in/~publ/ln/tifr47.pdf}.

\bibitem{KnusEtAl-BookOfInvols}
Knus M.A., Merkurjev A., Rost M., Tignol J.P., The book of involutions,
  \textit{American Mathematical Society Colloquium Publications}, Vol.~44,
  Amer. Math. Soc., Providence, RI, 1998.

\bibitem{Lam-IntroQuadForms}
Lam T.Y., Introduction to quadratic forms over f\/ields, \textit{Graduate Studies
  in Mathematics}, Vol.~67, Amer. Math. Soc., Providence, RI, 2005.

\bibitem{Morris-RRepsQForms}
Morris D., Real representations of semisimple {L}ie algebras have {${\mathbb
  Q}$}-forms, in Algebraic Groups and Arithmetic, Tata Institute of Fundamental
  Research, Mumbai, 2004, 469--490.

\bibitem{PlatonovRapinchukBook}
Platonov V., Rapinchuk A., Algebraic groups and number theory, \textit{Pure and
  Applied Mathematics}, Vol.~139, Academic Press, Inc., Boston, MA, 1994.

\bibitem{PrasadRapinchuk-Prescribed}
Prasad G., Rapinchuk A., On the existence of isotropic forms of semi-simple
  algebraic groups over number f\/ields with prescribed local behavior,
  \href{http://dx.doi.org/10.1016/j.aim.2006.01.001}{\textit{Adv. Math.}} \textbf{207} (2006), 646--660.

\bibitem{Raghunathan-ArithLattsSSGrps}
Raghunathan M.S., Arithmetic lattices in semisimple groups, \href{http://dx.doi.org/10.1007/BF02967980}{\textit{Proc.
  Indian Acad. Sci. Math. Sci.}} \textbf{91} (1982), 133--138.

\bibitem{Tits-Classification}
Tits J., Classif\/ication of algebraic semisimple groups, in Algebraic {G}roups
  and {D}iscontinuous {S}ubgroups ({P}roc. {S}ympos. {P}ure {M}ath., {B}oulder,
  {C}olo., 1965), Amer. Math. Soc., Providence, R.I., 1966, 33--62.

\bibitem{Tits-RepLinIrred}
Tits J., Repr\'esentations lin\'eaires irr\'eductibles d'un groupe r\'eductif
  sur un corps quelconque, \href{http://dx.doi.org/10.1515/crll.1971.247.196}{\textit{J.~Reine Angew. Math.}} \textbf{247} (1971),
  196--220.

\bibitem{Tits-StronglyInner}
Tits J., Strongly inner anisotropic forms of simple algebraic groups,
  \href{http://dx.doi.org/10.1016/0021-8693(90)90201-X}{\textit{J.~Algebra}} \textbf{131} (1990), 648--677.

\end{thebibliography}
\end{document}
